\documentclass{article}

% NeurIPS-style formatting
\usepackage[preprint]{neurips_2024}
\usepackage[utf8]{inputenc}
\usepackage[T1]{fontenc}
\usepackage{hyperref}
\usepackage{url}
\usepackage{booktabs}
\usepackage{amsfonts}
\usepackage{amsmath}
\usepackage{nicefrac}
\usepackage{microtype}
\usepackage{xcolor}
\usepackage{graphicx}
\usepackage{subcaption}

\title{Computational Graph Structure Distinguishes Surface Recall\\
from Explanatory Reasoning in Language Models}

\author{%
  darkflobi \AND flobi \\
  Independent Researchers \\
  \texttt{heyzoos123@gmail.com}
}

\begin{document}

\maketitle

\begin{abstract}
We investigate whether the distinction between surface-level factual recall and explanatory 
reasoning is detectable in the internal computation of language models, using sparse 
autoencoder (SAE) circuit attribution to analyze computational graph structure. Our study 
proceeds in three phases. First, we demonstrate a negative result: surface-level text 
statistics (TF-IDF, sequence length, N-gram overlap, readability) reliably classify surface 
vs.\ insight text, but this reflects vocabulary and length confounds rather than genuine 
reasoning differences. After five rounds of iterative confound control across three model 
scales (GPT-2 Small, Gemma-2-2B, Mistral 7B), we produce a vocabulary-controlled dataset 
($n$=36 pairs, TF-IDF = 50.1\%, length accuracy = 52.9\%, both at chance) on which a linear 
probe over GPT-2 residual stream activations finds no signal (50--56\% across all layers). 
Second, we pivot to direct circuit analysis using SAE attribution on Gemma-2-2B. In an 
exploratory analysis, surface-recall prompts concentrate feature activation at output-proximal 
layers (higher mean activated layer), while matched insight prompts engage more distributed 
mid-layer computation. Third, we pre-register the directional hypothesis 
($\Delta = \text{mean\_layer}(\text{surface}) - \text{mean\_layer}(\text{insight}) > 0$) 
and conduct two independent replications on held-out topic sets. A pre-registered replication 
($n$=11 valid of 15 total) yields $t = 2.04$, $p = 0.035$ (one-tailed, paired $t$-test), 
directionally consistent but not leave-one-out robust. A subsequent power replication on 30 
new held-out pairs yields $n$=30 valid (zero exclusions), $\Delta = {+}0.257$, 
$t(29) = 4.50$, $p = 0.000051$, Cohen's $d = 0.82$, 83.3\% directional, 
and full leave-one-out robustness. A prediction-confidence confound test finds 
$r = 0.001$ between confidence and depth deltas ($p = 0.996$), ruling out difficulty-driven 
explanations. We argue that the methodology---iterative confound elimination, 
pre-registered directional testing, sensitivity analysis, and confound controls---constitutes 
a contribution independent of the specific finding, and that computational graph depth 
offers a mechanistically grounded alternative to surface-level evaluation metrics.
\end{abstract}

\section{Introduction}

Evaluating whether a language model is reasoning or merely retrieving is among the hardest 
open problems in AI evaluation. Existing benchmarks measure what a model says---whether 
an answer is correct, fluent, or factually grounded---but cannot directly assess how the 
model arrives at that answer. A model that pattern-matches to a memorized response and a 
model that constructs a novel explanation may produce identical outputs while engaging 
fundamentally different internal computations.

This paper asks a simple question: can we detect this difference mechanistically? 
Specifically, can the structure of a language model's computational graph---which features 
activate, at which layers, in what distribution---distinguish prompts that elicit surface recall 
from prompts that elicit explanatory reasoning?

We begin with what does not work. Surface-level text statistics can classify surface vs.\ 
insight text with high accuracy, but this accuracy persists even when vocabulary is controlled, 
revealing it as a confound rather than a signal. Neural probes trained on residual stream 
activations find nothing once these confounds are eliminated.

We then turn to a different approach: direct analysis of computational graphs via sparse 
autoencoder (SAE) circuit attribution. Using the circuit-tracer framework \citep{lindsey2025} 
on Gemma-2-2B, we extract the set of SAE features activated during processing of each 
prompt and compute their mean layer depth. We find that surface-recall prompts concentrate 
activation at higher (output-proximal) layers, while insight prompts distribute activation more 
broadly across middle layers.

This finding survives pre-registered replication on held-out topics ($n$=30, $p < 0.001$, 
LOO robust), while surface statistics and linear probes fail once distributional confounds are 
controlled. A new confound test rules out prediction difficulty as an explanation 
($r = 0.001$, $p = 0.996$).

\subsection{Contributions}

\begin{enumerate}
\item \textbf{A methodological contribution:} We document a five-iteration confound 
elimination process showing how surface statistics systematically contaminate 
text-classification paradigms, with a reusable protocol for controlling them.

\item \textbf{A negative result:} Linear probes over GPT-2 residual stream activations cannot 
distinguish surface from insight text once distributional confounds are removed ($n$=36 
vocabulary-controlled pairs, all layers at chance). This constrains what probing studies can 
claim about reasoning representation.

\item \textbf{A positive result:} SAE circuit attribution reveals a robust difference in 
computational graph structure between surface recall and explanatory reasoning, replicated 
with pre-registration at $p < 0.001$ on held-out stimuli, with prediction confidence ruled 
out as a confound ($r = 0.001$).
\end{enumerate}

\section{Related Work}

\subsection{Mechanistic Interpretability and Sparse Autoencoders}

Recent work in mechanistic interpretability has moved beyond individual neuron analysis 
toward understanding circuits---connected subgraphs of computation implementing specific 
behaviors. Sparse autoencoders (SAEs) trained on residual stream activations recover 
interpretable monosemantic features distributed across layers 
\citep{cunningham2023, templeton2024}. Circuit tracing methods allow attribution of model 
outputs to specific active features, providing a graph-structured view of which features 
participate in each computation \citep{lindsey2025}.

\subsection{Reasoning Depth in Language Models}

Prior work on reasoning in language models has largely focused on behavioral benchmarks
---whether models produce correct answers to multi-step problems 
\citep{wei2022, kojima2022}. Less work characterizes the internal structure of that 
reasoning. Chain-of-thought prompting improves performance on reasoning tasks but does 
not characterize where in the network the reasoning occurs. Our work takes a different 
approach: rather than measuring output quality, we measure the depth profile of feature 
activation for prompts that vary in their reasoning demands while controlling for surface form.

\subsection{Layer Specialization in Transformers}

There is evidence that transformer layers are not uniform---earlier layers tend to encode 
syntactic and positional information, while later layers encode more semantic and task-specific 
content \citep{tenney2019, rogers2020}. If surface recall is more akin to pattern-matching on 
stored associations, we would expect it to engage later, more directly output-predictive layers. 
If insight requires integration of distributed knowledge, we would expect it to involve more 
mid-network processing. Our metric---mean layer of activated SAE features---operationalizes 
this intuition.

\section{Methods}

\subsection{Prompt Pair Construction}

We construct matched prompt pairs on the following criteria:
\begin{enumerate}
\item \textbf{Vocabulary match:} Surface and insight prompts for a given topic draw from the 
same lexical domain.
\item \textbf{Token count match:} Prompts are matched within $\pm$2 tokens.
\item \textbf{Syntactic structure match:} Both prompts are sentence-completion fragments of 
comparable syntactic form.
\item \textbf{Semantic contrast:} Surface prompts invoke direct factual statements; insight 
prompts invoke causal or mechanistic explanation of the same phenomenon.
\end{enumerate}

Example pair (topic: \textit{fire}):
\begin{itemize}
\item Surface: \textit{``What does fire need to burn?''} $\rightarrow$ target: \textit{oxygen}
\item Insight: \textit{``Why does fire spread once it starts?''} $\rightarrow$ target: \textit{because}
\end{itemize}

Both prompts concern fire, use similar vocabulary, and complete at the same syntactic 
position. The surface prompt invokes factual recall; the insight prompt requires encoding a 
mechanistic feedback relationship.

Pairs were constructed for 75+ topics across three phases: exploratory, held-out replication, 
and power replication. Topics were non-overlapping across phases.

\subsection{Negative Result: Surface Statistics}

Prior to circuit tracing, we conducted five rounds of statistical analysis on surface text 
properties:
\begin{itemize}
\item Zipf word frequency (per-token average log frequency)
\item Flesch-Kincaid readability grade level
\item Type-token ratio
\item Mean N-gram log probability (unigram, bigram, trigram)
\item TF-IDF classification accuracy
\item Sequence length accuracy
\end{itemize}

Across all metrics and all pairs, no consistent separation was found once vocabulary was 
controlled. The v4 dataset (word-set-enforced generation) achieved TF-IDF = 50.1\% and 
length accuracy = 52.9\%, both at chance.

\subsection{Circuit Tracing and Feature Depth}

We use \texttt{circuit\_tracer} \citep{lindsey2025} with Gemma-2-2B loaded via 
TransformerLens. For each prompt, we compute feature attribution graphs using the 
\texttt{attribute()} function with the following parameters:
\begin{itemize}
\item \texttt{batch\_size=32}
\item \texttt{max\_feature\_nodes=500}
\item \texttt{offload="cpu"}
\item \texttt{dtype=torch.bfloat16}
\end{itemize}

Hardware: RTX 4070 8GB GPU; attribution time $\approx$22s per prompt.

For each attribution graph, we extract the \texttt{active\_features} tensor of shape 
$(n_\text{active}, 3)$ where column 0 is the layer index. We compute the mean layer index 
across all active features as our depth metric:
\begin{equation}
\Delta = \overline{L}_\text{surface} - \overline{L}_\text{insight}
\end{equation}
where $\overline{L}$ denotes mean activated layer index. Positive $\Delta$ indicates surface 
prompts activate features at higher (output-proximal) layers.

\subsection{Pre-registration}

Before running the held-out replication, the directional hypothesis ($\Delta > 0$) was 
pre-registered in \texttt{experiment\_12\_replication\_pairs.json} via commit timestamp in 
the public repository. The statistical test (one-tailed paired $t$-test in the predicted 
direction) was specified in advance.

\subsection{Sensitivity Analysis}

Given the small effective sample size in the held-out replication ($n$=11 after exclusions), 
we ran leave-one-out (LOO) sensitivity analysis: for each pair, we removed it and 
recomputed the $t$-test on the remaining pairs. This tests whether any single pair is 
responsible for the significant result.

\section{Results}

\subsection{Surface Statistics Fail to Generalize Past Confound Control}

\begin{table}[h]
\centering
\caption{Surface statistics across all confound-control rounds.}
\label{tab:surface_stats}
\begin{tabular}{llrrr}
\toprule
Round & Method & $n$ pairs & TF-IDF & Length Acc. \\
\midrule
v1 & Unconstrained generation & $\sim$70 & 96.0\% & 73.0\% \\
v2 & Seed-sentence constraints & 24 & 74.4\% & 64.7\% \\
v3 & Surface-first rewrite & 6 clean & 74.4\% & --- \\
v4 & Word-set enforced & 36 & \textbf{50.1\%} & \textbf{52.9\%} \\
\bottomrule
\end{tabular}
\end{table}

This trajectory demonstrates that explanation is partially but not entirely constituted by 
vocabulary at the token level. The v4 generator enforced zero new words in insight rewrites, 
achieving confound control at 90\% yield rate (36/40 seeds).

\subsection{Neural Probes Find No Signal on Controlled Data}

On the v4 vocabulary-controlled dataset ($n$=36 pairs, TF-IDF at chance), logistic regression 
probes trained on GPT-2 (117M) residual stream activations achieve 50--56\% accuracy across 
all 12 layers, with no layer reaching statistical significance.

\begin{table}[h]
\centering
\caption{Linear probe accuracy by layer, GPT-2 Small, vocabulary-controlled pairs.}
\label{tab:probe}
\begin{tabular}{lr@{\hspace{1.5em}}r}
\toprule
Layer & Accuracy & $p$-value \\
\midrule
0 & 52.8\% & 0.41 \\
3 & 55.6\% & 0.31 \\
6 & 50.0\% & 0.50 \\
9 & 52.8\% & 0.44 \\
11 & 50.0\% & 0.50 \\
\bottomrule
\end{tabular}
\end{table}

This null result rules out the concern that our subsequent positive finding is driven by 
residual vocabulary confounds. If distributional confounds were present, the probe would 
detect them---as it did in rounds v1--v3.

\subsection{Exploratory Circuit Attribution}

Using SAE circuit attribution on Gemma-2-2B, we compute mean activated layer depth for 
each prompt in the exploratory set (30 pairs submitted; 2 valid after CUDA attribution failures). 
Both valid pairs showed surface depth greater than insight depth (100\%), mean 
$\Delta = {+}0.326$.

This finding reversed our initial hypothesis: we originally predicted that insight prompts 
would activate higher layers (equating ``deeper reasoning'' with ``deeper layers''). The data 
showed the opposite, consistent with a direct-lookup vs.\ distributed-computation 
interpretation. The revised hypothesis ($\Delta > 0$, surface $>$ insight) was formulated 
after seeing the exploratory data and was pre-registered before running any replication data.

\subsection{Pre-registered Replication ($n$=11)}

Fifteen held-out pairs were constructed on new topics. Four were excluded due to 
attribution extraction failure (multi-token target tokens). Results on 11 valid pairs:

\begin{table}[h]
\centering
\caption{Pre-registered replication results ($n$=11 valid pairs).}
\label{tab:rep1}
\begin{tabular}{lr}
\toprule
Metric & Value \\
\midrule
Surface mean layer & 8.378 \\
Insight mean layer & 8.260 \\
$\Delta$ (surface $-$ insight) & $+$0.118 \\
Pairs in predicted direction & 8/11 (72.7\%) \\
Paired $t$-test (one-tailed) & $t = 2.04$, $p = 0.035$ \\
Wilcoxon signed-rank (one-tailed) & $p = 0.051$ \\
\bottomrule
\end{tabular}
\end{table}

\textbf{LOO sensitivity:} The result is not LOO-robust: removing any one of 7/11 pairs 
pushes the parametric $p$-value above 0.05 (range: 0.050--0.070). The directional finding 
is stable across all subsets; the fragility reflects statistical power at $n$=11, not an 
outlier-driven result. This motivated the power replication.

\subsection{Power Replication ($n$=30)}

Thirty new held-out pairs were constructed on topics with zero overlap with all 45 prior pairs. 
All 60 target tokens were pre-audited for single-token status (18 initially multi-token targets 
were replaced). The same pre-registered hypothesis and tests were applied.

\begin{table}[h]
\centering
\caption{Power replication results ($n$=30, all valid, zero exclusions).}
\label{tab:rep2}
\begin{tabular}{lr}
\toprule
Metric & Value \\
\midrule
Pairs submitted / valid & 30 / 30 \\
Surface mean layer & 8.359 \\
Insight mean layer & 8.102 \\
$\Delta$ (surface $-$ insight) & $+$0.257 \\
Cohen's $d$ & 0.82 \\
Pairs in predicted direction & 25/30 (83.3\%) \\
Paired $t$-test (one-tailed) & $t(29) = 4.500$, $p = 0.000051$ \\
Wilcoxon signed-rank (one-tailed) & $p < 0.001$ \\
LOO robust & Yes (all 30 subsets $p < 0.05$) \\
\bottomrule
\end{tabular}
\end{table}

The power replication confirms the directional finding with strong statistical evidence. Both 
parametric and non-parametric tests converge at $p < 0.001$. The effect size ($+$0.257) is 
consistent with the exploratory estimate ($+$0.326), suggesting the exploratory estimate was 
accurate and the initial replication ($+$0.118) was attenuated by low power rather than 
reflecting the true effect size.

\subsection{Prediction Confidence Confound Test}

A natural alternative explanation is that surface-recall prompts are simply more predictable 
(lower perplexity), and that prediction difficulty---not semantic category---drives the depth 
effect. To test this, we computed $\log P(\text{target} \mid \text{prompt})$ for all 30 power 
pairs using Gemma-2-2B and tested whether prediction confidence differences correlate with 
depth differences.

\begin{table}[h]
\centering
\caption{Prediction confidence confound test ($n$=30 power pairs).}
\label{tab:confound}
\begin{tabular}{lr}
\toprule
Metric & Value \\
\midrule
Surface more confident ($\Delta_\text{lp} > 0$) & 60.0\% of pairs \\
Mean $\Delta_\text{lp}$ (surface $-$ insight) & $+$1.97 (SD = 5.21) \\
$t$-test vs.\ 0 & $t = 2.04$, $p = 0.051$ \\
Pearson $r(\Delta_\text{lp},\, \Delta_\text{depth})$ & $r = 0.001$, $p = 0.996$ \\
Spearman $\rho$ & $\rho = -0.008$, $p = 0.966$ \\
\bottomrule
\end{tabular}
\end{table}

Surface prompts are marginally more confident on average (60\% of pairs, $p = 0.051$), 
which is consistent with factual completions being more predictable. However, the correlation 
between confidence differences and depth differences is essentially zero ($r = 0.001$, 
$p = 0.996$). This means the pairs where surface prompts are most confident are no more 
likely to show large depth differentials than pairs where surface and insight prompts have 
similar confidence. \textbf{The depth effect is semantically driven, not difficulty-driven.}

\section{Discussion}

\subsection{Summary of Findings}

Surface-level text statistics and linear probes fail to distinguish surface recall from 
explanatory reasoning once distributional confounds are controlled. However, SAE circuit 
attribution reveals a consistent structural difference: surface-recall prompts concentrate 
feature activation at output-proximal layers, while insight prompts engage more distributed 
mid-layer computation. This difference is robust to pre-registered replication on held-out 
topics ($p < 0.001$), survives leave-one-out sensitivity analysis, and is not explained by 
prediction confidence differences ($r = 0.001$).

\subsection{On HARKing and Directional Reversal}

The exploratory finding reversed our initial hypothesis. We originally predicted that insight 
prompts would activate higher layers (equating ``deeper reasoning'' with ``deeper layers''). 
The data showed the opposite: surface prompts activate higher layers, consistent with a 
direct-lookup interpretation where factual retrieval engages output-proximal features while 
explanatory reasoning requires distributed mid-layer composition.

We treat this reversal transparently. The revised hypothesis was formulated after seeing 
the exploratory data and was pre-registered before running any replication data. 
The replication design---new topics, pre-registered direction, held-out data---prevents this 
from constituting HARKing \citep{kerr1998}.

\subsection{Exclusion Rate Trajectory}

The exclusion rate across phases (0\% exploratory $\to$ 27\% initial replication $\to$ 0\% 
power replication) reflects pipeline maturation, not selective analysis. The excluded 
replication pairs failed because target tokens were multi-token under the Gemma-2-2B 
tokenizer; the power replication eliminated this by pre-auditing all targets. The excluded 
pairs span diverse domains and have no systematic relationship to the hypothesis.

\subsection{Limitations}

\textbf{Single model.} All positive results are on Gemma-2-2B. Whether the mean-layer-depth 
metric generalizes to other architectures, model sizes, or training regimes is unknown.

\textbf{Controlled stimuli.} The matched-pair paradigm constrains both topics and prompt 
structure. Whether the depth metric discriminates on unconstrained, naturalistic text is 
untested.

\textbf{Single metric.} Mean activated layer depth is a coarse summary. Richer metrics 
(layer variance, feature co-activation patterns, graph diameter) might capture additional 
dimensions.

\textbf{Sample size.} While the power replication ($n$=30) provides strong statistical 
evidence, 30 pairs across a single model is a small foundation for broad claims.

\subsection{The Methodology as Contribution}

Independent of the specific finding, this paper documents a complete methodology for 
mechanistic evaluation of text-processing differences: iterative confound control, vocabulary 
enforcement, baseline validation, pre-registered directional testing with non-parametric 
robustness checks, leave-one-out sensitivity analysis, and prediction-confidence confound 
control. This protocol is applicable to any study attempting to detect representational 
differences between text categories in language model activations.

\subsection{Implications for Evaluation}

If the mean-layer-depth metric generalizes, it offers a fundamentally different evaluation 
paradigm: measuring \textit{how} a model processes a prompt, not \textit{what} it produces. 
This is harder to game than output-based metrics---a model cannot change its computational 
graph structure through prompt engineering or output formatting. The depth score reflects 
the internal architecture of computation, providing a ground-truth signal about processing 
depth that is inaccessible from the output alone.

\section{Conclusion}

We find that computational graph structure, measured via SAE circuit attribution, 
distinguishes surface recall from explanatory reasoning in Gemma-2-2B. Surface-recall 
prompts concentrate feature activation at output-proximal layers; explanatory-reasoning 
prompts distribute activation across middle layers. This difference survives pre-registered 
replication on held-out topics ($n$=30, $p < 0.001$, LOO robust, Cohen's $d = 0.82$), 
while surface statistics and linear probes fail once distributional confounds are controlled, 
and prediction confidence is ruled out as a confound ($r = 0.001$). Whether this metric 
generalizes across models, scales, and text distributions remains an open question and a 
clear direction for future work.

Code, data, and pre-registration records:
\url{https://github.com/heyzoos123-blip/sae-reasoning-depth}

\section*{Acknowledgments}

We thank the authors of \texttt{circuit\_tracer} for releasing their tool publicly, enabling 
this line of research.

\bibliographystyle{plainnat}
\begin{thebibliography}{10}

\bibitem[Cunningham et al.(2023)]{cunningham2023}
Cunningham, H., Ewert, A., Riggs, L., Huben, R., \& Sharkey, L. (2023).
\newblock Sparse autoencoders find highly interpretable features in language models.
\newblock \textit{arXiv preprint arXiv:2309.08600}.

\bibitem[Kerr(1998)]{kerr1998}
Kerr, N.~L. (1998).
\newblock HARKing: Hypothesizing after the results are known.
\newblock \textit{Personality and Social Psychology Review}, 2(3), 196--217.

\bibitem[Kojima et al.(2022)]{kojima2022}
Kojima, T., Gu, S.~S., Reid, M., Matsuo, Y., \& Iwasawa, Y. (2022).
\newblock Large language models are zero-shot reasoners.
\newblock \textit{Advances in Neural Information Processing Systems}, 35.

\bibitem[Lindsey et al.(2025)]{lindsey2025}
Lindsey, J., et al. (2025).
\newblock Circuit-tracer: Feature attribution via sparse autoencoder circuit tracing.
\newblock \textit{Anthropic Technical Report}.

\bibitem[Nanda et al.(2023)]{nanda2023}
Nanda, N., Chan, L., Lieberum, T., Smith, J., \& Steinhardt, J. (2023).
\newblock Progress measures for grokking via mechanistic interpretability.
\newblock \textit{International Conference on Learning Representations}.

\bibitem[Rogers et al.(2020)]{rogers2020}
Rogers, A., Kovaleva, O., \& Rumshisky, A. (2020).
\newblock A primer in BERTology: What we know about how BERT works.
\newblock \textit{Transactions of the Association for Computational Linguistics}, 8, 842--866.

\bibitem[Templeton et al.(2024)]{templeton2024}
Templeton, A., et al. (2024).
\newblock Scaling monosemanticity: Extracting interpretable features from Claude 3 Sonnet.
\newblock \textit{Anthropic Research Blog}.

\bibitem[Tenney et al.(2019)]{tenney2019}
Tenney, I., Das, D., \& Pavlick, E. (2019).
\newblock BERT rediscovers the classical NLP pipeline.
\newblock \textit{Proceedings of ACL}, 4593--4601.

\bibitem[Wei et al.(2022)]{wei2022}
Wei, J., Wang, X., Schuurmans, D., Bosma, M., Xia, F., Chi, E., Le, Q.~V., \& Zhou, D. (2022).
\newblock Chain-of-thought prompting elicits reasoning in large language models.
\newblock \textit{Advances in Neural Information Processing Systems}, 35.

\end{thebibliography}

\appendix

\section{Surface Statistics Across All Rounds}

\begin{table}[h]
\centering
\caption{Complete surface statistics by dataset version.}
\begin{tabular}{llrrrl}
\toprule
Round & Method & $n$ & TF-IDF & Length Acc. & Notes \\
\midrule
v1 & Unconstrained & $\sim$70 & 96.0\% & 73.0\% & Confounded \\
v2 & Seed-sentence & 24 & 74.4\% & 64.7\% & Partially controlled \\
v3 & Surface-first rewrite & 6 & 74.4\% & --- & Low yield \\
v4 & Word-set enforced & 36 & 50.1\% & 52.9\% & At chance \\
\bottomrule
\end{tabular}
\end{table}

\section{Sign Convention and Experiment Log}

\begin{table}[h]
\centering
\caption{Sign conventions across experiments.}
\begin{tabular}{lll}
\toprule
Experiment & JSON convention & Paper convention \\
\midrule
exp11 (exploratory) & $\Delta = \text{insight} - \text{surface}$ & Sign flipped \\
exp12 (replication) & $\Delta = \text{surface} - \text{insight}$ & Unchanged \\
Power replication & $\Delta = \text{surface} - \text{insight}$ & Unchanged \\
\bottomrule
\end{tabular}
\end{table}

\section{Excluded Pairs (Replication Phase)}

\begin{table}[h]
\centering
\caption{Pairs excluded from pre-registered replication.}
\begin{tabular}{lll}
\toprule
Topic & Failure Mode & Direction if scorable \\
\midrule
interest\_rates & 0 active features (surface) & Unknown \\
democracy\_decay & 0 active features (insight) & Unknown \\
reputation & 0 active features (surface) & Unknown \\
momentum\_social & 0 active features (both) & Unknown \\
\bottomrule
\end{tabular}
\end{table}

All four exclusions are attributable to target tokens resolving to multiple BPE tokens under 
the Gemma-2-2B tokenizer, causing the circuit-tracer to fail feature extraction. No excluded 
pair is scorable in either direction. Exclusion is technical, not directional.

\section{Leave-One-Out Analysis (Replication, $n$=11)}

LOO analysis shows 7/11 removals push $p > 0.05$ (range: 0.050--0.070); 4/11 retain 
significance. The three pairs whose removal improves the result are those with negative or 
near-zero deltas---consistent with distributional prediction, not evidence of outlier dependence. 
The power replication achieves full LOO robustness.

\section{Power Replication LOO Summary}

All 30 single-pair removals retain $p < 0.05$. Minimum $p$ across all LOO iterations: 
$p < 0.001$. The result is fully robust to any single-pair removal.

\end{document}
